\subsection{Class-three exclusions after every finite original prefix}
\label{sec:ghost-dense-class3}

Let \(\mathscr A\) be the set of infinite prime sequences
allowed by the original A.44 recursion, with exponents
\(n_j=(j+4)!\). A finite prefix consists of complete
pairs \((q_i,r_i)\), together with legal finite CRT
residue choices through those stages. Write \(\operatorname{pr}_k\) for
restriction to the first \(k\) pairs. Give this space
its prefix topology, in which fixing finitely many
of these choices specifies a basic open set.
Let \(\mathscr A_c\) be the subspace whose support
primes satisfy every compact elliptic order test (16).
Finally let \(\mathscr F_3\) be the subspace whose
first empty unrestricted integral nilpotent section
fibre is class three.

\begin{theorem}[Dense class-three failure with the original factorial exponents]
\label{thm:ghost-dense-class3}
For every \(k\ge0\), the actual prefix maps satisfy
\[
 \begin{split}
 \operatorname{pr}_k(\mathscr F_3)
       &=\operatorname{pr}_k(\mathscr A),\\
 \operatorname{pr}_k(\mathscr F_3\cap\mathscr A_c)
       &=\operatorname{pr}_k(\mathscr A_c).
 \end{split}                                         \tag{64}
\]
Thus class-three-excluded outputs are dense among
all original A.44 outputs, and among the outputs
passing the compact tests. In the latter domain,
every constructed pair has a full profinite section
of the affine-cover abelian quotient with proper
class zero, while its metabelian cover section
fibre is empty.

More precisely, after any specified finite prefix
the construction chooses a coefficient prime
\(\ell\ge5\) and a later support prime
\(q_{\ell-4}\) at which every integral class-two
section has nonzero third obstruction, already
detected modulo \(\ell^2\). All exponents remain
the original \((j+4)!\).
\end{theorem}
\begin{proof}
Fix a legal prefix of length \(k\). Choose a prime
\(\ell>k+4\), \(\ell\ge5\), outside the finite
set of its support primes, and put
\(J=\ell-4\). Fill the intervening stages through
\(J-1\) by the original rules, excluding the one
prime \(\ell\) from the new choices. There are
infinitely many eligible primes at each step, so
this finite exclusion is allowed. In the compact-
passing domain, impose additionally the noncube
conditions of
Proposition~\ref{prop:ghost-elliptic-passing} at
every new stage. Their previously proved field
intersections preserve every legal CRT class.
Write \(a=a_J,b=b_J\) for the resulting integers,
and let \(S_0\) be the finite set of fresh support
primes preceding stage \(J\).
Then
\[
 \begin{gathered}
 n=n_J=\ell!,\qquad m=n_{J+1}=(\ell+1)!,\\
 v_\ell(n)=1,\qquad \ell\nmid n_i\ (i<J),\qquad
 \ell\nmid ab,\\
 v_3(a)=v_2(b)=1,\qquad v_2(a)=v_3(b)=0.
 \end{gathered}                                      \tag{65}
\]

Choose the stage-\(J\) local roots \(\xi_p,\eta_p\)
required by A.44. At \(\ell\) take
\(\xi_\ell=1\). The original invariant and
Hensel estimate read
\[
 \begin{split}
 v_\ell(a+b-1)&\ge H_J(\ell)=\ell+1,\\
 v_\ell(\eta_\ell-1)&\ge\ell,\qquad
       v_\ell(M)=H_{J+1}(\ell)=\ell+2,
 \end{split}                                         \tag{66}
\]
where \(M\) is the original stage modulus.
In particular its CRT class \(Q_0\) has
\(Q_0\equiv1\pmod{\ell^2}\), and still has
\(Q_0\equiv1\pmod m\).

We give the coefficient-\(\ell\) version of the
prefix comparison with all earlier support primes
retained. Choose positive integer placeholders
\(P,R\), prime to \(S_0\), agreeing with
\(\xi_p,\eta_p\) modulo \(p\) for each
\(p\in S_0\). For fixed integral Kummer cocycles
\(\alpha=\kappa_\ell(a)\),
\(\beta=\kappa_\ell(b)\), set
\[
 a^*=\alpha+n\kappa_\ell(P),\qquad
 b^*=\beta+n\kappa_\ell(R).                            \tag{67}
\]
At every old positive support prime, the rational
number \(bR^n\) is a principal unit; at every old
negative support prime, \(aP^n\) is a principal
unit. Their integral local Kummer primitives
therefore give exactly the central inertia values
\(n_i u_p^*\) and \(-n_i v_p^*\), respectively,
reduced modulo \(\ell\).

The finite cup of \((a,b)\) modulo \(\ell\)
vanishes: the original invariant supplies an
ordinary local point of \(aX^n+bY^n=1\) at every
prime, and \(\ell\mid n\). Finite global detection
and the residue surjection of
Proposition~\ref{prop:ghost-ordinary-cup} supply
a global central primitive \(h_*\) with exactly
the preceding old inertia values and zero inertia
outside \(S_0\cup\{2,3,\ell\}\).
The uniqueness group is
\[
 H^1(G_{\mathbb Q,\{\ell,\infty\}},
                      \mathbb F_\ell(2))=0.          \tag{68}
\]
Indeed its integral first cohomology injects into
\(H^1(G_{\mathbb Q},\mathbb Z_\ell(2))=0\),
by Proposition~\ref{prop:ghost-h1-weight2}; its
integral second cohomology is
\(K_2(\mathbb Z)\{\ell\}=0\), by the same
checked comparison used for (52). The coefficient
sequence proves (68). At two and three the scalar
\(p^2-p\) is invertible modulo \(\ell\), so no
extra inertia ambiguity remains there.

These prescriptions freeze the reduction of every
integral class-two section for every continuation
with
\[
 q_i\equiv r_i\equiv1\pmod{\prod_{p\in S_0}p}
                            \quad(i>J).              \tag{69}
\]
Here is the exact cochain comparison. Choose its
two literal integral Kummer cocycles as the
uniformly convergent sums
\[
 \begin{split}
 a_\infty&=\alpha+n\kappa_\ell(q_J)
                         +\sum_{i>J}n_i\kappa_\ell(q_i),\\
 b_\infty&=\beta+n\kappa_\ell(r_J)
                         +\sum_{i>J}n_i\kappa_\ell(r_i).
 \end{split}
\]
At an old prime, \(q_J/P,r_J/R\) and all later
prime factors in (69) are principal units. Divide
the differences from (67) by \(\ell\). Every
summand is an integral local coboundary, and the
coefficients \(n_i/\ell\) tend to zero.
The integral primitive series converges and gives
\[
 a_\infty-a^*=\ell\delta e_p,
 \qquad b_\infty-b^*=\ell\delta f_p
                       \quad(p\in S_0).              \tag{70}
\]
Thus each old local primitive agrees modulo
\(\ell\) with the prescribed one. At every new
support prime, \(\ell\mid n_i\), so the actual
integral tame relation makes central inertia zero
modulo \(\ell\). At zero-valuation primes it is
already zero integrally. Hence the difference
between any reduced integral lift and \(h_*\)
is unramified outside \(\ell\); (68) makes it
a central coboundary. This proves the comparison
with no finite-coefficient division.

The finite primitive just constructed defines
\(\rho_*:G_{\mathbb Q}\to
H_\ell\rtimes\mathbb F_\ell^\times\), with fixed
field \(L_*\). Valuations at three and two give
independent abelian Kummer characters over
\(\mathbb Q(\mu_\ell)\). Commutators therefore
give the whole Heisenberg group \(H_\ell\), of
order \(\ell^3\). Since multiplication by
\(2-1\) on the abelianization is invertible,
the commutator subgroup of the full image is
\(H_\ell\). Consequently
\[
 \operatorname{Gal}(L_*/\mathbb Q)
       =H_\ell\rtimes\mathbb F_\ell^\times,
       \qquad L_*^{\rm ab}=\mathbb Q(\mu_\ell).
                                                        \tag{71}
\]

For the actual Chebotarev comparison put
\[
 F=\mathbb Q(\mu_M),\quad
 R_b=F(b^{1/n}),\quad
 \mathcal R=F(b^{1/n},a^{1/\ell}).                     \tag{72}
\]
The original A.44 argument extends the CRT
automorphism \(\sigma_F\) to \(R_b\), fixing
\(b^{1/n}\). Rational \(\ell\)-power classes
inject into a cyclotomic field: a nontrivial
class has a prime valuation not divisible by
\(\ell\), so its radical degree is \(\ell\);
a normal radical extension of that degree would
contain \(\mu_\ell\), whose degree \(\ell-1\)
does not divide \(\ell\).
The valuation independence of \(a,b\) therefore
persists over \(F\). Since \(v_\ell(n)=1\),
the unique degree-\(\ell\) subfield of the cyclic
extension \(R_b/F\) is \(F(b^{1/\ell})\),
disjoint from \(F(a^{1/\ell})\).
Thus \(\sigma_F\) extends to \(\mathcal R\),
fixing both indicated radicals. Moreover (71)
and abelianity of \(\mathcal R/F\) give the
exact intersection
\[
 L_*F\cap\mathcal R=F(a^{1/\ell},b^{1/\ell}),
               \qquad [L_*\mathcal R:\mathcal R]=\ell.
                                                        \tag{73}
\]
A nonzero central element of \(H_\ell\), paired
with \(\sigma_F\), fixes both roots in (73).
It agrees there with the specified automorphism
of \(\mathcal R\), and hence extends to the
actual compositum. This verifies compatibility
with the possibly nontrivial cyclotomic CRT action.

For the compact-passing construction also adjoin
\(C=F(3^{1/3})\). The unique cubic subfield of
\(R_b/F\) is \(F(b^{1/3})\), independent of
\(C\) by valuations at two and three and (22).
Adjoining \(a^{1/\ell}\), and then the extension
of degree \(\ell\) in (73), introduces no cubic
intersection, since \(\ell\ge5\). Thus
\(C\cap L_*\mathcal R=F\), and the action
\(3^{1/3}\mapsto\zeta_3 3^{1/3}\), with the
same \(\sigma_F\), extends to that compositum.
Chebotarev now gives a prime \(q_J\) satisfying
every original condition and
\[
 q_J\equiv1\pmod{\ell^2},\qquad
 \rho_*(\operatorname{Frob}_{q_J})=(0,0,c;1),
                 \quad c\in\mathbb F_\ell^\times.
                                                        \tag{74}
\]
In the compact-passing construction it additionally
has three a noncube. Choose \(r_J\) by the original
rules, with two a noncube in that construction.

Condition (69) remains available at every later
stage. The stage-\(J\) CRT equations have made
the coefficient sum one modulo every old support
prime. If such a prime does not divide the current
exponent, the unit-coefficient variable has a
simple root near one. If it divides that exponent,
the original invariant gives the stronger precision
\(H_i(p)>2v_p(n_i)\), so the original Hensel
estimate again gives a root one modulo \(p\).
The additional congruence required when
\(p\mid n_{i+1}\) is retained by the same bound.
Taking these roots gives (69) by CRT; every other
original condition and every extra noncube choice
remains available. This constructs a legal infinite
continuation without changing its prescribed prefix.

For any integral class-two section of that completed
pair, (70) and (68) identify its reduction with
\(\rho_*\). At \(q=q_J\), (74) therefore gives
\(\alpha\equiv0\), \(\gamma\equiv c\ne0\)
modulo \(\ell\), unaffected by a central
coboundary. With \(m_0=n_J=\ell!\), its actual
third residue is
\[
 Z_{D,q}=m_0(\gamma-q\alpha u)
             +(q^2-1)u\frac{m_0(m_0-1)}2.
                                                        \tag{75}
\]
It has \(\ell\)-adic valuation one, while
\(q^2-1\) has valuation at least two. Hence every
integral lift is obstructed, already modulo
\(\ell^2\). Class-two existence for every
ordinary pair proves that this first failure is
exactly class three. In the compact-passing
domain, Proposition~\ref{prop:ghost-all-passing-genjac}
gives the affine abelian cover section, and the
actual morphism (63) excludes every metabelian
cover lift. This proves both equalities in (64).
\end{proof}

The exact symbolic calculation in
\path{code/check_dense_class3.py} starts from both
full Hall residues of
Proposition~\ref{prop:ghost-class3-residue}, verifies
(75), and proves the integral polynomial identity
\(2Z_{D,q}\equiv2\ell(m_0/\ell)\gamma\pmod{\ell^2}\)
under its stated congruences. It also checks the
original factorial and Hensel indices at twelve
coefficient primes from five through forty-three;
the proof of (65)--(66) applies to every prime
\(\ell\ge5\).

The map from a prime sequence to its two completed
products is continuous: at a fixed power modulus
all sufficiently late factorial exponents are
zero, so that modulus depends on a finite prefix.
Consequently the same exclusions are dense in the
corresponding sets of completed pairs for their
induced power-completion topology. Density proves
the surjectivities (64); it makes no assertion that
every infinite continuation is excluded. The
universal all-depth question and the unspecified
remaining compact-passing continuations retain
their exact unresolved status.
